\section{Syntaktická analýza a LL(1) gramatika}

\subsection{Princíp}
Syntaktická analýza je založená na rekurzívnom zostupe podľa LL(1) gramatiky uvedenej nižšie (modul \texttt{parser.c}). Každý neterminál je implementovaný jednou funkciou, ktorá podľa aktuálneho tokenu vyberie pravidlo, a tokeny sa spotrebúvajú funkciou \texttt{compare\_token}. Tokeny parser neberie priamo zo skenera, ale z poľa, do ktorého ich pri inicializácii načíta všetky. Výrazy sa nespracúvajú rekurzívnym zostupom, ale samostatným precedenčným analyzátorom (modul \texttt{expr.c}, podsekcia \ref{sec:expr}), ktorý parser volá všade, kde gramatika obsahuje neterminál $\langle\text{EXP}\rangle$.

Parser nielen overuje syntax, ale súčasne \textbf{buduje abstraktný syntaktický strom} (\texttt{ast.h}) a vykonáva sémantické kontroly (kapitola o sémantickej analýze). Každá funkcia neterminálu vracia ukazovateľ na vytvorený uzol cez výstupný parameter a vlastníctvo uzla prechádza na volajúceho. Pri chybe funkcia uvoľní všetko, čo si dovtedy alokovala, takže sa pri žiadnom chybovom výstupe neuniká pamäť.

\subsection{Abstraktný syntaktický strom}
Koreňom stromu je uzol \texttt{NODE\_PROGRAM} so zoznamom funkcií a rozšírení (\texttt{extension}). Uzly sa delia na:
\begin{itemize}
    \item \textbf{definície:} \texttt{NODE\_FUNC\_DEF}, \texttt{NODE\_EXTENSION}, \texttt{NODE\_PARAM}, \texttt{NODE\_BLOCK},
    \item \textbf{príkazy:} \texttt{NODE\_VAL\_DEF}, \texttt{NODE\_VAR\_DEF}, \texttt{NODE\_ASSIGN}, \texttt{NODE\_IF\_STMT}, \texttt{NODE\_WHILE}, \texttt{NODE\_FOR}, \texttt{NODE\_RETURN}, \texttt{NODE\_EXPR\_STMT},
    \item \textbf{výrazy:} literály, \texttt{NODE\_ID}, \texttt{NODE\_UNARY\_OP}, \texttt{NODE\_BINARY\_OP}, \texttt{NODE\_CALL}, \texttt{NODE\_MEMBER}, \texttt{NODE\_ISTRING} a \texttt{NODE\_IF\_EXPR}.
\end{itemize}
Každý uzol nesie dátový typ (\texttt{DataType}) určený pri jeho vytvorení a pozíciu v zdrojovom kóde. Funkcie a volania s viacerými zoznamami parametrov (\texttt{f(a)(b, c)}) ukladajú parametre a argumenty ``naplocho'' spolu s poľom \texttt{list\_sizes}, ktoré hovorí, koľko z nich patrí do ktorých zátvoriek. Telo funkcie zapísané jedným výrazom sa uloží ako blok s jediným uzlom \texttt{NODE\_RETURN}.

\subsection{Predbežný prechod a dopredné volania}
Aby sa dala funkcia volať pred svojou definíciou, prejde parser po načítaní tokenov celé pole ešte raz (funkcia \texttt{prepass}). Sleduje pritom odsadenie a hľadá kľúčové slovo \texttt{def} na najvyššej úrovni a v telách rozšírení. Z každej hlavičky prečíta názov, zoznamy parametrov a prípadný návratový typ a zaregistruje ju v globálnej tabuľke symbolov spolu s pozíciou definície v poli tokenov. Hlavičku, ktorej nerozumie, preskočí, a chybu ohlási až samotný parser, keď na ňu dôjde.

Samotné parsovanie ide potom ako predtým rekurzívnym zostupom. Pri dosiahnutí \texttt{def} už existujúci symbol iba doplní, a ak je už označený ako definovaný, ide o redefiníciu. Návratový typ, ktorý nie je uvedený, sa odvodí z tela funkcie. Ak ho však potrebuje volanie skôr, než sa k telu dostane hlavný prechod, funkcia \texttt{resolve\_function} telo naparsuje vopred (\texttt{parse\_ahead}): uloží stav parsera, presunie sa na tokeny danej definície, vytvorí jej uzol a stav vráti späť. Hotový uzol sa zapamätá a keď hlavný prechod na túto definíciu príde, uzol len prevezme a tokeny preskočí. Poradie funkcií v AST tak zostáva rovnaké ako v zdrojovom texte. Ak sa funkcia bez uvedeného typu volá z vlastného tela (priamo alebo cez inú funkciu), jej typ ešte nie je známy a ide o chybu 6.

\subsection{Rozsahy platnosti}
Parser udržiava zreťazené tabuľky symbolov (\texttt{current\_scope}). Nový rozsah vzniká pre telo funkcie (spolu s parametrami), každý odsadený blok, jednoriadkový príkaz vo vetve \texttt{if}/\texttt{else} a v tele cyklu a pre cyklus \texttt{for} (riadiaca premenná je viditeľná iba v~strážach a tele). Premenná vzniká v tabuľke až po spracovaní svojho inicializačného výrazu, takže \texttt{val x = x + 1} sa správne odmietne ako použitie nedefinovaného identifikátora.

\subsection{LL(1) gramatika}
\small
\noindent
\begin{longtable}{p{0.5cm} p{3.0cm} p{6.2cm} p{3.8cm}}
\toprule
\textbf{Č.} & \textbf{Neterminál} & \textbf{Pravidlo (pravá strana)} & \textbf{Stručný popis} \\
\midrule
\endhead

\bottomrule
\endfoot

1 & $\langle\text{PROLOG}\rangle$ & $\rightarrow$ import id . * EOL $\langle\text{F\_LIST}\rangle$ EOF & Začiatok programu; identifikátor musí byť \texttt{ifj26}. \\
2 & $\langle\text{F\_LIST}\rangle$ & $\rightarrow$ $\langle\text{F\_DEF}\rangle$ $\langle\text{F\_LIST}\rangle$ & Ďalšia definícia funkcie. \\
3 & $\langle\text{F\_LIST}\rangle$ & $\rightarrow$ $\langle\text{EXTENSION}\rangle$ $\langle\text{F\_LIST}\rangle$ & Ďalšie rozšírenie typu. \\
4 & $\langle\text{F\_LIST}\rangle$ & $\rightarrow$ $\varepsilon$ & Koniec zoznamu definícií. \\
5 & $\langle\text{F\_DEF}\rangle$ & $\rightarrow$ $\langle\text{ANNOTATION}\rangle$ def id $\langle\text{PARAM\_LISTS}\rangle$ $\langle\text{D\_TYPE}\rangle$ = $\langle\text{F\_BODY\_END}\rangle$ & Definícia funkcie s voliteľnou anotáciou. \\
6 & $\langle\text{PARAM\_LISTS}\rangle$ & $\rightarrow$ ( $\langle\text{PARAMETER}\rangle$ ) $\langle\text{PARAM\_LISTS}\rangle$ & Zoznam parametrov; môže ich byť viac (currying). \\
7 & $\langle\text{PARAM\_LISTS}\rangle$ & $\rightarrow$ $\varepsilon$ & Koniec zoznamov parametrov. \\
8 & $\langle\text{F\_BODY\_END}\rangle$ & $\rightarrow$ EOL* INDENT $\langle\text{F\_BODY}\rangle$ DEDENT & Viacriadkové telo funkcie. \\
9 & $\langle\text{F\_BODY\_END}\rangle$ & $\rightarrow$ $\langle\text{EXP}\rangle$ EOL & Telo tvorené jedným výrazom (implicitný \texttt{return}). \\
10 & $\langle\text{ANNOTATION}\rangle$ & $\rightarrow$ @ main $\langle\text{NEW\_LINE}\rangle$ & Označenie funkcie \texttt{@main}. \\
11 & $\langle\text{ANNOTATION}\rangle$ & $\rightarrow$ $\varepsilon$ & Funkcia bez anotácie. \\
12 & $\langle\text{NEW\_LINE}\rangle$ & $\rightarrow$ EOL & Nový riadok za anotáciou. \\
13 & $\langle\text{NEW\_LINE}\rangle$ & $\rightarrow$ $\varepsilon$ & Bez nového riadku. \\
14 & $\langle\text{EXTENSION}\rangle$ & $\rightarrow$ extension ( id : $\langle\text{TYPE}\rangle$ ) EOL INDENT $\langle\text{E\_BODY}\rangle$ DEDENT & Rozšírenie typu o metódy; \texttt{id} je prijímač. \\
15 & $\langle\text{E\_BODY}\rangle$ & $\rightarrow$ def id $\langle\text{PARAM\_LISTS}\rangle$ $\langle\text{D\_TYPE}\rangle$ = $\langle\text{F\_BODY\_END}\rangle$ $\langle\text{EB\_NEXT}\rangle$ & Metóda rozšírenia. \\
16 & $\langle\text{EB\_NEXT}\rangle$ & $\rightarrow$ $\langle\text{E\_BODY}\rangle$ & Ďalšia metóda. \\
17 & $\langle\text{EB\_NEXT}\rangle$ & $\rightarrow$ $\varepsilon$ & Koniec metód. \\
18 & $\langle\text{D\_TYPE}\rangle$ & $\rightarrow$ : $\langle\text{TYPE}\rangle$ & Explicitný typ. \\
19 & $\langle\text{D\_TYPE}\rangle$ & $\rightarrow$ $\varepsilon$ & Typ nie je uvedený (odvodí sa). \\
20--24 & $\langle\text{TYPE}\rangle$ & $\rightarrow$ Int $\mid$ Double $\mid$ String $\mid$ Unit $\mid$ Boolean & Dátové typy. \\
25 & $\langle\text{PARAMETER}\rangle$ & $\rightarrow$ id : $\langle\text{TYPE}\rangle$ $\langle\text{PARAM\_NEXT}\rangle$ & Prvý parameter. \\
26 & $\langle\text{PARAMETER}\rangle$ & $\rightarrow$ $\varepsilon$ & Prázdny zoznam parametrov. \\
27 & $\langle\text{PARAM\_NEXT}\rangle$ & $\rightarrow$ , $\langle\text{PARAM\_TAIL}\rangle$ & Pokračovanie zoznamu. \\
28 & $\langle\text{PARAM\_NEXT}\rangle$ & $\rightarrow$ $\varepsilon$ & Koniec zoznamu. \\
29 & $\langle\text{PARAM\_TAIL}\rangle$ & $\rightarrow$ id : $\langle\text{TYPE}\rangle$ $\langle\text{PARAM\_NEXT}\rangle$ & Ďalší parameter. \\
30 & $\langle\text{PARAM\_TAIL}\rangle$ & $\rightarrow$ $\varepsilon$ & Koniec zoznamu (koncová čiarka). \\
31 & $\langle\text{F\_BODY}\rangle$ & $\rightarrow$ $\langle\text{STMT}\rangle$ $\langle\text{STMT\_NEXT}\rangle$ & Telo bloku. \\
32 & $\langle\text{STMT\_NEXT}\rangle$ & $\rightarrow$ EOL $\langle\text{STMT}\rangle$ $\langle\text{STMT\_NEXT}\rangle$ & Príkazy oddelené koncom riadku. \\
33 & $\langle\text{STMT\_NEXT}\rangle$ & $\rightarrow$ ; $\langle\text{STMT}\rangle$ $\langle\text{STMT\_NEXT}\rangle$ & Príkazy oddelené bodkočiarkou. \\
34 & $\langle\text{STMT\_NEXT}\rangle$ & $\rightarrow$ $\varepsilon$ & Koniec postupnosti príkazov; príkaz smie nasledovať priamo za blokom bez oddeľovača. \\
35 & $\langle\text{STMT}\rangle$ & $\rightarrow$ val id $\langle\text{D\_TYPE}\rangle$ = $\langle\text{EXP}\rangle$ & Nemenná premenná. \\
36 & $\langle\text{STMT}\rangle$ & $\rightarrow$ var id $\langle\text{D\_TYPE}\rangle$ = $\langle\text{EXP}\rangle$ & Menná premenná. \\
37 & $\langle\text{STMT}\rangle$ & $\rightarrow$ while $\langle\text{EXP}\rangle$ do EOL INDENT $\langle\text{F\_BODY}\rangle$ DEDENT & Cyklus \texttt{while}. \\
38 & $\langle\text{STMT}\rangle$ & $\rightarrow$ return $\langle\text{RET\_VAL}\rangle$ & Návrat z funkcie. \\
39 & $\langle\text{STMT}\rangle$ & $\rightarrow$ $\langle\text{EXP}\rangle$ & Výraz ako príkaz (napr. volanie funkcie). \\
40 & $\langle\text{STMT}\rangle$ & $\rightarrow$ $\langle\text{EXP}\rangle$ = $\langle\text{EXP}\rangle$ & Priradenie; ľavá strana musí byť identifikátor. \\
41 & $\langle\text{RET\_VAL}\rangle$ & $\rightarrow$ $\langle\text{EXP}\rangle$ & Návrat s hodnotou. \\
42 & $\langle\text{RET\_VAL}\rangle$ & $\rightarrow$ $\varepsilon$ & Návrat bez hodnoty. \\
48 & $\langle\text{STMT}\rangle$ & $\rightarrow$ $\langle\text{IF\_STMT}\rangle$ & Podmienený príkaz. \\
49 & $\langle\text{IF\_STMT}\rangle$ & $\rightarrow$ if $\langle\text{EXP}\rangle$ [then] ( EOL INDENT $\langle\text{F\_BODY}\rangle$ DEDENT $\mid$ $\langle\text{STMT}\rangle$ ) $\langle\text{IF\_NEXT}\rangle$ & Vetva \texttt{then} ako blok alebo jeden príkaz. \\
50 & $\langle\text{IF\_NEXT}\rangle$ & $\rightarrow$ else $\langle\text{ELSE\_NEXT}\rangle$ & Vetva \texttt{else}. \\
51 & $\langle\text{IF\_NEXT}\rangle$ & $\rightarrow$ $\varepsilon$ & Bez \texttt{else}. \\
52 & $\langle\text{ELSE\_NEXT}\rangle$ & $\rightarrow$ $\langle\text{IF\_STMT}\rangle$ & Reťazenie \texttt{else if}. \\
53 & $\langle\text{ELSE\_NEXT}\rangle$ & $\rightarrow$ EOL INDENT $\langle\text{F\_BODY}\rangle$ DEDENT $\mid$ $\langle\text{STMT}\rangle$ & Blok alebo jeden príkaz. \\
54 & $\langle\text{IF\_CONTROL}\rangle$ & $\rightarrow$ if $\langle\text{EXP}\rangle$ [then] $\langle\text{IF\_CTRL\_NEXT}\rangle$ else $\langle\text{IF\_CTRL\_NEXT}\rangle$ & \texttt{if} ako výraz; \texttt{else} je povinný. \\
55 & $\langle\text{IF\_CTRL\_NEXT}\rangle$ & $\rightarrow$ $\langle\text{EXP}\rangle$ & Vetva tvorená výrazom. \\
56 & $\langle\text{IF\_CTRL\_NEXT}\rangle$ & $\rightarrow$ EOL INDENT $\langle\text{F\_BODY}\rangle$ DEDENT & Vetva tvorená blokom. \\
57 & $\langle\text{STMT}\rangle$ & $\rightarrow$ for $\langle\text{FOR\_CONTENT}\rangle$ & Cyklus \texttt{for}. \\
59 & $\langle\text{FOR\_CONTENT}\rangle$ & $\rightarrow$ $\langle\text{GENERATOR}\rangle$ $\langle\text{OPT\_GUARDS}\rangle$ do $\langle\text{FOR\_BODY}\rangle$ & Jednoriadková hlavička. \\
60 & $\langle\text{FOR\_CONTENT}\rangle$ & $\rightarrow$ EOL INDENT $\langle\text{GENERATOR}\rangle$ $\langle\text{EOL\_GUARDS}\rangle$ DEDENT do $\langle\text{FOR\_BODY}\rangle$ & Viacriadková hlavička. \\
61 & $\langle\text{GENERATOR}\rangle$ & $\rightarrow$ id <- $\langle\text{EXP}\rangle$ $\langle\text{RANGE}\rangle$ $\langle\text{EXP}\rangle$ $\langle\text{OPT\_STEP}\rangle$ & Generátor hodnôt. \\
62, 63 & $\langle\text{RANGE}\rangle$ & $\rightarrow$ to $\mid$ until & Horná hranica vrátane / bez nej. \\
64 & $\langle\text{OPT\_STEP}\rangle$ & $\rightarrow$ by $\langle\text{EXP}\rangle$ & Krok. \\
65 & $\langle\text{OPT\_STEP}\rangle$ & $\rightarrow$ $\varepsilon$ & Bez kroku. \\
66 & $\langle\text{OPT\_GUARDS}\rangle$ & $\rightarrow$ if $\langle\text{EXP}\rangle$ $\langle\text{OPT\_GUARDS}\rangle$ & Stráž cyklu. \\
67 & $\langle\text{OPT\_GUARDS}\rangle$ & $\rightarrow$ $\varepsilon$ & Bez stráží. \\
68 & $\langle\text{EOL\_GUARDS}\rangle$ & $\rightarrow$ EOL if $\langle\text{EXP}\rangle$ $\langle\text{EOL\_GUARDS}\rangle$ & Stráž na samostatnom riadku. \\
69 & $\langle\text{EOL\_GUARDS}\rangle$ & $\rightarrow$ $\varepsilon$ & Bez ďalších stráží. \\
70 & $\langle\text{FOR\_BODY}\rangle$ & $\rightarrow$ $\langle\text{STMT}\rangle$ & Telo cyklu jedným príkazom. \\
71 & $\langle\text{FOR\_BODY}\rangle$ & $\rightarrow$ EOL INDENT $\langle\text{F\_BODY}\rangle$ DEDENT & Telo cyklu blokom. \\
\end{longtable}
\normalsize

\noindent Čísla pravidiel zodpovedajú komentárom v zdrojovom kóde \texttt{parser.c}. Pravidlá 43--47 pôvodnej verzie (neterminály \texttt{VAL}, \texttt{DEC}, \texttt{ARG}, \texttt{TERM}) zanikli, pretože ich nahradil precedenčný analyzátor výrazov, ktorý spracúva aj volania funkcií.

\subsection{Spracovanie výrazov (precedenčná syntaktická analýza)}
\label{sec:expr}
Výrazy spracúva modul \texttt{expr.c} pomocou zdola nahor riadeného precedenčného analyzátora so zásobníkom a precedenčnou tabuľkou. Zásobník obsahuje terminály (operátory, zátvorky, koncová značka \$) a neterminály $E$, ku ktorým je pripojený podstrom AST a jeho typ.

\begin{table}[!ht]
\centering
\small
\begin{tabular}{l|ccccccccccc}
 & \texttt{!} & $-_u$ & $*\;/$ & $+\;-$ & rel & $=\;\neq$ & \&\& & $\mid\mid$ & ( & ) & \$ \\
\hline
\texttt{!} & $<$ & $<$ & $>$ & $>$ & $>$ & $>$ & $>$ & $>$ & $<$ & $>$ & $>$ \\
$-_u$ & $<$ & $<$ & $>$ & $>$ & $>$ & $>$ & $>$ & $>$ & $<$ & $>$ & $>$ \\
$*\;/$ & $<$ & $<$ & $>$ & $>$ & $>$ & $>$ & $>$ & $>$ & $<$ & $>$ & $>$ \\
$+\;-$ & $<$ & $<$ & $<$ & $>$ & $>$ & $>$ & $>$ & $>$ & $<$ & $>$ & $>$ \\
rel & $<$ & $<$ & $<$ & $<$ & -- & $>$ & $>$ & $>$ & $<$ & $>$ & $>$ \\
$=\;\neq$ & $<$ & $<$ & $<$ & $<$ & $<$ & -- & $>$ & $>$ & $<$ & $>$ & $>$ \\
\&\& & $<$ & $<$ & $<$ & $<$ & $<$ & $<$ & $>$ & $>$ & $<$ & $>$ & $>$ \\
$\mid\mid$ & $<$ & $<$ & $<$ & $<$ & $<$ & $<$ & $<$ & $>$ & $<$ & $>$ & $>$ \\
( & $<$ & $<$ & $<$ & $<$ & $<$ & $<$ & $<$ & $<$ & $<$ & $=$ & -- \\
) & -- & -- & -- & -- & -- & -- & -- & -- & -- & -- & -- \\
\$ & $<$ & $<$ & $<$ & $<$ & $<$ & $<$ & $<$ & $<$ & $<$ & -- & acc \\
\end{tabular}
\caption{Precedenčná tabuľka tak, ako je v poli \texttt{PREC} v \texttt{expr.c} (riadok = terminál na vrchole zásobníka, stĺpec = vstupný terminál; rel = \texttt{< <= > >=}). Značka $<$ je posun, $>$ redukcia, $=$ spojenie zátvoriek, \texttt{--} chyba a \emph{acc} koniec výrazu.}
\end{table}

Priorita od najvyššej: unárne \texttt{!} a \texttt{-}, \texttt{*} \texttt{/}, \texttt{+} \texttt{-}, relačné operátory, \texttt{==} \texttt{!=}, \texttt{\&\&}, \texttt{||}. Operátory \texttt{+ - * /}, \texttt{\&\&} a \texttt{||} sú ľavo asociatívne. Relačné operátory (\texttt{< <= > >=}) a operátory rovnosti (\texttt{== !=}) asociatívne nie sú: v tabuľke je pre dvojicu rel--rel a rovnosť--rovnosť chyba, takže napríklad \texttt{a < b < c} je syntaktická chyba (kód 2). Ďalšie vlastnosti:
\begin{itemize}
    \item \textbf{Unárne a binárne mínus} sa rozlišujú podľa predchádzajúceho symbolu: ak bol predtým operand, ide o binárny operátor, inak o unárny.
    \item \textbf{Operandy} (literály, identifikátory, volania funkcií, \texttt{if} výraz a interpolované reťazce) sa spracujú rekurzívnymi pomocnými funkciami a uložia sa na zásobník ako neterminál $E$. Identifikátor funkcie sa spracuje ako volanie vrátane viacerých zoznamov argumentov \texttt{f(a)(b, c)}.
    \item \textbf{Prístup k členu} \texttt{x.meno} a \texttt{x.meno(args)} sa viaže silnejšie ako každý operátor. Hľadá sa v tomto poradí: metóda rozšírenia pre typ prijímača, funkcia s daným menom, ktorá berie prijímač ako prvý argument, a nakoniec vstavaný člen \texttt{toString}, \texttt{toInt}, \texttt{floor} alebo \texttt{length} (iba bez zátvoriek a len na vhodnom type).
    \item \textbf{Prázdne zátvorky} \texttt{()} tvoria hodnotu typu \texttt{Unit}.
    \item \textbf{Koniec výrazu} sa rozpozná na prvom tokene, ktorý ho nemôže predĺžiť (EOL, \texttt{then}, \texttt{do}, čiarka, nepárová \texttt{)} a pod.); analyzátor ho nespotrebuje.
    \item Pri každej redukcii $E \rightarrow E\ op\ E$ resp. $E \rightarrow op\ E$ sa kontrolujú typy operandov, určí sa typ výsledku a vytvorí sa uzol \texttt{NODE\_BINARY\_OP} resp. \texttt{NODE\_UNARY\_OP}.
\end{itemize}

\medskip
\noindent Samotný analyzátor funguje takto. Vstupným bodom \texttt{expr.c} je funkcia \texttt{parse\_expr}, ktorá pracuje so zásobníkom položiek \texttt{PItem}, kde každá položka je buď terminál (operátor, zátvorka, koncová značka), alebo neterminál $E$ s podstromom AST a jeho typom. Na začiatku je na zásobníku iba značka \$ a príznak \texttt{prev\_operand} hovorí, či bol posledný spracovaný symbol operand. V cykle sa pre aktuálny token urobí tento postup:
\begin{enumerate}
    \item Ak za operandom nasleduje bodka, zavolá sa \texttt{parse\_member} a operand na vrchole zásobníka sa nahradí uzlom \texttt{NODE\_MEMBER}.
    \item Ak token začína operand a predtým operand nebol, \texttt{parse\_primary} z neho vytvorí uzol a uloží ho na zásobník ako $E$. Podľa tokenu je to literál, interpolovaný reťazec (\texttt{interp\_string}), \texttt{if} výraz (\texttt{if\_control} z parsera) alebo identifikátor (\texttt{parse\_identifier}). Identifikátor premennej dostane typ zo symbolu. Pri funkcii sa načítajú zoznamy argumentov (\texttt{parse\_arg\_lists}), zavolá sa \texttt{resolve\_function}, \texttt{validate\_function\_call} skontroluje počet, štruktúru zoznamov a typy a vznikne uzol \texttt{NODE\_CALL}.
    \item Dvojica \texttt{()} sa nahradí hodnotou \texttt{Unit}.
    \item Inak sa token funkciou \texttt{classify} priradí k terminálu tabuľky. Token, ktorý výraz nemôže predĺžiť, sa zmení na koncovú značku \$.
    \item Potom sa porovná terminál najbližšie k vrcholu zásobníka so vstupným terminálom v tabuľke \texttt{PREC}. Pri posune sa vstupný terminál uloží a token sa spotrebuje. Pri redukcii sa zavolá \texttt{reduce}, ktorá odoberie operátor a operandy, skontroluje typy (\texttt{check\_unary}, \texttt{check\_binary}), vytvorí uzol a vloží späť výsledné $E$. Pri \texttt{=} sa zátvorky odstránia. Chyba v tabuľke ukončí analýzu kódom 2, nevhodné typy operandov pri redukcii kódom 6 a pri \emph{acc} je výraz hotový.
\end{enumerate}
Po prijatí musí na zásobníku zostať práve značka \$ a jediné $E$, ktorého uzol a typ sa vrátia volajúcemu. Zásobník vlastní všetky ešte nezačlenené podstromy, takže pri chybe v ktoromkoľvek kroku sa uvoľnia a nič neunikne.
